\chapter{Reproducing the paper}
\label{ch:repro}

Before changing anything, I ran the published code and data of \citet{bian2025} on our lab computer
(NVIDIA RTX A2000, 12~GB) to check that I can reach their accuracy.

\section{Setup}
The authors' dataset contains 3,893 training and 1,126 test samples. Each sample holds the \cotwo{}
at 7,492 points of the room for the last 12 time steps (6~minutes), the 13 control values (number of
people and the airflow rate and angle of 6 vent groups) and the \cotwo{} of the next 6 time steps
(3~minutes). The network (GNOT, 552,969 parameters) was trained for 200 epochs, which took about
6~hours per network. I trained five networks with the seeds 2023 to 2027 and used their mean as
the forecast, as in the paper.

\section{Result}
\tref{tab:repro} compares my networks with Table~3 of the paper (relative $\ell_2$ error).

\begin{table}[h]
  \centering
  \caption{Reproduction of \citet{bian2025}: relative $\ell_2$ error of the single networks and of
  the ensemble.}
  \label{tab:repro}
  \begin{tabular}{lrrrr}
    \toprule
     & \multicolumn{2}{c}{This work} & \multicolumn{2}{c}{Paper}\\
     & Train & Test & Train & Test\\
    \midrule
    Network 1 & 5.85\,\% & 12.70\,\% & 6.35\,\% & 12.09\,\%\\
    Network 2 & 5.88\,\% & 12.71\,\% & 6.33\,\% & 11.83\,\%\\
    Network 3 & 5.84\,\% & 12.70\,\% & 6.33\,\% & 11.82\,\%\\
    Network 4 & 5.92\,\% & 12.53\,\% & 6.35\,\% & 12.74\,\%\\
    Network 5 & 5.81\,\% & 12.97\,\% & 6.33\,\% & 13.01\,\%\\
    Ensemble  & 5.48\,\% & 11.46\,\% & 5.9\,\%  & 10.90\,\%\\
    \bottomrule
  \end{tabular}
\end{table}

The single networks are in the same range as in the paper (12.5--13.0\,\% against 11.8--13.0\,\%),
and the ensemble is clearly better than every single network, as in the paper. My ensemble error is
0.56 percentage points higher than the published one. One difference is the loss: the training
script in the repository trains networks that predict only the mean \cotwo{} with a relative
$\ell_2$ loss, while the paper's networks also predict a variance and are trained with a Gaussian
negative log-likelihood. Different random seeds and hardware may also contribute.

\section{Problems in the published code}
To run the code I had to fix a few things:
\begin{itemize}
  \item the training script called the likelihood loss with the wrong arguments and stopped with
        an error; I used the relative $\ell_2$ loss of the repository instead;
  \item the variant with mean and variance output (GNOTE) failed because a tuple was treated as a
        tensor, and its likelihood loss was never used;
  \item a data file needed by the control notebook had been deleted in a later commit and had to be
        taken from an older commit;
  \item the notebooks expected a second GPU (\texttt{cuda:1});
  \item the code only runs with specific package versions (PyTorch 2.3.0, DGL 2.2.1, NumPy below 2.0).
\end{itemize}
I packed the working version with a single evaluation command, so that it can be run again from
the repository \todo{repository link}.
